% =====================================================================
\appendix
\section{Phụ lục}
\label{sec:appendix}
% =====================================================================

\subsection{Kiểm tra kỹ thuật tự động}
\label{sec:appendix-selftest}

Trước khi huấn luyện, script \code{scripts/self\_test.py} kiểm tra bốn nhóm thuộc tính và ghi
kết quả vào \code{results/self\_test.json}. Bảng dưới đây tóm tắt các kiểm tra và tiêu chí đạt.

\begin{table}[htbp]
  \centering
  \small
  \begin{tabular}{@{}p{5.2cm} p{5.0cm} l@{}}
    \toprule
    Kiểm tra & Tiêu chí & Kết quả \\
    \midrule
    Chuỗi kích thước tensor & Khớp bảng \ref{tab:arch} &
    \textcolor{m8green}{\textbf{ĐẠT}} \\
    Số tham số suy luận & $11.169.858$ & \textcolor{m8green}{\textbf{ĐẠT}} \\
    Số tham số huấn luyện & $11.171.910$ & \textcolor{m8green}{\textbf{ĐẠT}} \\
    Forward/backward ba cấu hình & Không lỗi, logits đúng chiều & \textcolor{m8green}{\textbf{ĐẠT}} \\
    \code{torch.rot90} khớp nhãn $k$ & Sai số $< 10^{-5}$ & \textcolor{m8green}{\textbf{ĐẠT}} \\
    Không trùng split & $\Dtrain \cap \Dval = \varnothing$ & \textcolor{m8green}{\textbf{ĐẠT}} \\
    Tập nhãn lồng nhau & $\DL^{10} \subset \DL^{20} \subset \DL^{50}$ & \textcolor{m8green}{\textbf{ĐẠT}} \\
    Không rò rỉ nhãn & $\DL \cap \DU = \varnothing$ và $(\DL \cup \DU) \cap \Dval = \varnothing$
    & \textcolor{m8green}{\textbf{ĐẠT}} \\
    Số lượng theo \tabref{tab:label-budget} & Đúng $4.500 / 9.000 / 22.500$ &
    \textcolor{m8green}{\textbf{ĐẠT}} \\
    \bottomrule
  \end{tabular}
  \caption{Kết quả kiểm tra kỹ thuật tự động trên môi trường thực nghiệm. Script trả mã thoát
  khác $0$ nếu bất kỳ kiểm tra nào không đạt, nên quy trình có thể được dùng trong CI.}
  \label{tab:selftest}
\end{table}

\subsection{Gói mã nguồn bàn giao}
\label{sec:appendix-repro}

\paragraph{Cấu trúc thư mục.}
\begin{lstlisting}[language=bash,caption={Cấu trúc gói tái lập.},label={lst:tree}]
D:\IOT
|- .venv/                     # moi truong ao Python 3.12
|- requirements.txt
|- data/
|  |- cifar-10-batches-py/    # du lieu goc
|  |- artifacts/splits.json   # chi muc split/subset da dong bang
|- src/m8/
|  |- config.py               # hang so giao thuc (ke ca sieu tham so contrastive)
|  |- data.py                 # anh xa sieu lop, split, ngan sach nhan, view contrastive
|  |- models.py               # ResNet-18 32x32 + MultiTaskNet (3 head tuy chon)
|  |- contrastive.py          # Projection Head, NT-Xent, SupCon
|  |- train.py                # vong huan luyen E1..E6
|  |- evaluate.py             # Macro-F1, Accuracy, per-class, confusion
|  |- benchmark.py            # tham so, dung luong, do tre
|  |- experiments.py          # dieu phoi khoi chinh va khoi mo rong
|- scripts/
|  |- prepare_data.py         # tai du lieu + kiem tra split
|  |- self_test.py            # kiem tra ky thuat tu dong
|  |- check_contrastive.py    # kiem tra NT-Xent/SupCon bang vi du tinh tay
|  |- sweep_lambda_c.py       # khao sat lambda_c CHI tren validation
|  |- lock_lambda_c.py        # ap quy tac khoa lambda_c
|  |- run_experiments.py      # khoi thuc nghiem chinh (27 luot)
|  |- run_extensions.py       # khoi mo rong contrastive (E3--E6)
|  |- run_inference_benchmark.py
|  |- run_rotation_diagnostic.py   # Rotation Accuracy chan doan
|  |- run_rotation_data_cost.py    # chi phi dung view cua nhanh phu
|  |- make_dataset_figures.py # hinh minh hoa du lieu (muc 2)
|  |- make_figures.py         # sinh hinh cho bao cao
|  |- make_extension_figures.py    # hinh khoi mo rong
|  |- make_report_data.py     # sinh bang/macro LaTeX
|  |- run_error_analysis.py   # du doan test theo chi muc + logits (muc 5.2)
|  |- make_error_figure.py    # hinh cac truong hop du doan sai
|  |- package_delivery.py     # dong goi checkpoint + MANIFEST co SHA-256
|  |- verify_macros.py, verify_tables.py, verify_pdf.py, inventory.py
|- notebooks/
|  |- M8_CIFAR10_MultiTask.ipynb
|  |- M8_Final_All.ipynb      # notebook cuoi cung, tong hop toan bo de tai
|- results/
|  |- checkpoints/*.json|*.pt # ket qua va checkpoint khoi chinh
|  |- checkpoints_ext/*.json|*.pt  # ket qua khoi mo rong E3--E6
|  |- summary.json            # tong hop khoi chinh
|  |- summary_all.json        # tong hop gop CA HAI khoi (chi khac summary.json khi
|                             #   da co ket qua mo rong trong checkpoints_ext/)
|  |- lambda_c_sweep.json     # khao sat lambda_c (chi validation)
|  |- benchmark_audit.json    # ket qua kiem tra che nhan (ban sao tu data/artifacts/)
|  |- split_indices.json      # DANH SACH CHI MUC split/subset da dong bang (3 MB)
|  |- self_test.json, environment.json, inference_benchmark.json
|  |- rotation_diagnostic.json, rotation_data_cost.json
|  |- runs_table.csv, figures/*.png
|  |- test_predictions/*.npz  # du doan test theo chi muc (de kiem tra lai chi so)
|- report/                    # bao cao LaTeX nay
\end{lstlisting}

\begin{keybox}
\textbf{Về checkpoint và gói bàn giao.} Đề cương §5.2 yêu cầu đóng gói checkpoint đã chọn bằng
validation cho các lượt chính. Checkpoint nặng khoảng $46$~MB mỗi lượt (tổng $\sim1{,}4$~GB cho
$39$ lượt), nên \textbf{không} đưa thẳng vào kho mã nguồn; \code{.gitignore} loại
\code{results/checkpoints/}, \code{results/checkpoints\_ext/} và mọi tệp \code{*.pt}.

Để vẫn thoả yêu cầu, kho mã nguồn kèm sẵn:

\begin{itemize}[leftmargin=1.1cm,itemsep=2pt]
  \item \code{results/test\_predictions/<tag>.npz} — dự đoán test \emph{theo đúng chỉ mục} gốc
        $0$--$9999$, kèm \code{logits} và \code{prob} softmax. Mỗi tệp chỉ $\sim150$~KB nên
        \textbf{được bàn giao}, đúng yêu cầu \enquote{kèm logits hoặc dự đoán test theo chỉ mục}.
  \item \code{scripts/package\_delivery.py} — đóng gói checkpoint thành tệp riêng kèm
        \code{MANIFEST.json} có băm SHA-256 của từng tệp. Script \emph{kiểm chứng lại chỉ số}
        của mỗi checkpoint trước khi đóng gói: nạp lại, chạy trên $\Dtest$, so với chỉ số đã ghi
        trong \code{<tag>.json}; dung sai đo bằng số ảnh bị lật.
\end{itemize}

Bản thân \emph{chỉ số} của mọi lượt chạy vẫn được bàn giao đầy đủ dưới dạng
\code{results/runs\_table.csv} và \code{summary*.json}, nên mọi bảng trong báo cáo đều kiểm tra
lại được. Muốn dựng lại checkpoint, chạy \code{scripts/run\_experiments.py} và
\code{scripts/run\_extensions.py} như mục dưới; kết quả đã có sẽ được bỏ qua nên chỉ tốn thời
gian cho phần còn thiếu.
\end{keybox}

\paragraph{Quy trình chạy lại từ đầu.}
\begin{lstlisting}[language=bash,caption={Tái lập toàn bộ kết quả từ dữ liệu gốc.},label={lst:repro}]
cd D:\IOT
.venv\Scripts\python.exe scripts\prepare_data.py          # tai du lieu + dong bang split
.venv\Scripts\python.exe scripts\self_test.py             # kiem tra ky thuat
.venv\Scripts\python.exe scripts\run_experiments.py       # 27 luot khoi chinh
.venv\Scripts\python.exe scripts\run_inference_benchmark.py
.venv\Scripts\python.exe scripts\run_rotation_diagnostic.py  # Rotation Accuracy
.venv\Scripts\python.exe scripts\run_rotation_data_cost.py   # chi phi dung view
.venv\Scripts\python.exe scripts\make_figures.py          # sinh hinh khoi chinh
.venv\Scripts\python.exe scripts\run_error_analysis.py    # du doan test theo chi muc (logits)
.venv\Scripts\python.exe scripts\make_error_figure.py     # hinh cac ca du doan sai
.venv\Scripts\python.exe scripts\make_report_data.py      # sinh bang LaTeX
.venv\Scripts\python.exe scripts\package_delivery.py --out delivery_m8 --zip  # dong goi
cd report && tectonic -X compile main.tex                 # bien dich bao cao
\end{lstlisting}

\paragraph{Chạy khối mở rộng E3--E6.}
Khối mở rộng có thêm hai bước khóa siêu tham số \emph{chỉ trên validation}:
\begin{lstlisting}[language=bash,caption={Quy trình khối mở rộng, tách bạch thông tin với test.},label={lst:ext}]
.venv\Scripts\python.exe scripts\sweep_lambda_c.py --epochs 40   # CHI validation
.venv\Scripts\python.exe scripts\lock_lambda_c.py                # ap quy tac khoa
.venv\Scripts\python.exe scripts\run_extensions.py               # 12 luot E3--E6
.venv\Scripts\python.exe scripts\make_extension_figures.py       # hinh khoi mo rong
.venv\Scripts\python.exe scripts\make_report_data.py             # sinh lai bang + macro
\end{lstlisting}
Script khảo sát \textbf{không} gọi bước đánh giá test, nên kết quả test không thể ảnh hưởng
tới quyết định khóa $\lamc$.

\paragraph{Tính xác định và khả năng tái lập.}
Mọi nguồn ngẫu nhiên đều được ghim: seed chia split $2026$; seed lần lặp $42$, $43$, $44$ dùng
cho khởi tạo trọng số, thứ tự dữ liệu, augmentation và bộ sinh số riêng của nhánh phụ. Chỉ mục
đầy đủ của mọi tập được lưu ở \code{data/artifacts/splits.json}. Trên phần cứng có kiến trúc
GPU khác, kết quả số sẽ không trùng khít vì một số phép tính trên GPU không có tính kết hợp
chặt, nhưng quy trình và phép chia dữ liệu vẫn tái lập được~\cite{pytorch_reproducibility}.

\subsection{Chi tiết từng lượt chạy}
\label{sec:appendix-perrun}

% Sinh tự động — scripts/make_report_data.py
\begingroup
\scriptsize
\setlength{\tabcolsep}{4pt}
\begin{longtable}{@{}l l r r r r r r@{}}
  \caption{Chi tiết từng lượt chạy trong khối thực nghiệm chính. \enquote{Epoch tốt} là epoch được chọn theo Macro-F1 validation.}
  \label{tab:per-run} \\
  \toprule
  Cấu hình & Nhãn & Lần lặp & Epoch tốt & Val Macro-F1 & Test Macro-F1 & Test Acc & Giây \\
  \midrule
  \endfirsthead
  \multicolumn{8}{l}{\small\itshape (tiếp theo \tabref{tab:per-run})} \\
  \toprule
  Cấu hình & Nhãn & Lần lặp & Epoch tốt & Val Macro-F1 & Test Macro-F1 & Test Acc & Giây \\
  \midrule
  \endhead
  \bottomrule
  \endfoot
  \cfg{E1} & 10\,\% & 42 & 98 & 0,9394 & 0,9359 & 0,9385 & 340,2 \\
  \cfg{E1} & 10\,\% & 43 & 70 & 0,9345 & 0,9305 & 0,9337 & 342,4 \\
  \cfg{E1} & 10\,\% & 44 & 76 & 0,9294 & 0,9294 & 0,9324 & 349,7 \\
  \cfg{E2} & 10\,\% & 42 & 85 & 0,9546 & 0,9510 & 0,9530 & 578,2 \\
  \cfg{E2} & 10\,\% & 43 & 79 & 0,9485 & 0,9450 & 0,9473 & 578,7 \\
  \cfg{E2} & 10\,\% & 44 & 78 & 0,9432 & 0,9451 & 0,9474 & 579,3 \\
  \cfg{E2-L} & 10\,\% & 42 & 90 & 0,9401 & 0,9329 & 0,9355 & 576,0 \\
  \cfg{E2-L} & 10\,\% & 43 & 88 & 0,9322 & 0,9279 & 0,9306 & 574,4 \\
  \cfg{E2-L} & 10\,\% & 44 & 67 & 0,9328 & 0,9332 & 0,9360 & 576,2 \\
  \cfg{E1} & 20\,\% & 42 & 100 & 0,9549 & 0,9516 & 0,9535 & 553,2 \\
  \cfg{E1} & 20\,\% & 43 & 97 & 0,9512 & 0,9513 & 0,9533 & 568,0 \\
  \cfg{E1} & 20\,\% & 44 & 100 & 0,9471 & 0,9462 & 0,9484 & 552,5 \\
  \cfg{E2} & 20\,\% & 42 & 88 & 0,9663 & 0,9668 & 0,9681 & 1.014,2 \\
  \cfg{E2} & 20\,\% & 43 & 82 & 0,9669 & 0,9645 & 0,9659 & 1.014,0 \\
  \cfg{E2} & 20\,\% & 44 & 81 & 0,9628 & 0,9579 & 0,9594 & 1.013,3 \\
  \cfg{E2-L} & 20\,\% & 42 & 91 & 0,9516 & 0,9519 & 0,9537 & 1.007,4 \\
  \cfg{E2-L} & 20\,\% & 43 & 85 & 0,9544 & 0,9522 & 0,9540 & 1.027,1 \\
  \cfg{E2-L} & 20\,\% & 44 & 90 & 0,9497 & 0,9475 & 0,9496 & 1.046,0 \\
  \cfg{E1} & 50\,\% & 42 & 97 & 0,9704 & 0,9695 & 0,9707 & 1.213,0 \\
  \cfg{E1} & 50\,\% & 43 & 100 & 0,9716 & 0,9702 & 0,9714 & 1.141,4 \\
  \cfg{E1} & 50\,\% & 44 & 79 & 0,9704 & 0,9673 & 0,9686 & 1.141,9 \\
  \cfg{E2} & 50\,\% & 42 & 81 & 0,9794 & 0,9783 & 0,9792 & 2.324,5 \\
  \cfg{E2} & 50\,\% & 43 & 100 & 0,9827 & 0,9796 & 0,9804 & 2.320,4 \\
  \cfg{E2} & 50\,\% & 44 & 91 & 0,9756 & 0,9753 & 0,9762 & 2.322,1 \\
  \cfg{E2-L} & 50\,\% & 42 & 85 & 0,9761 & 0,9739 & 0,9749 & 2.330,9 \\
  \cfg{E2-L} & 50\,\% & 43 & 89 & 0,9779 & 0,9765 & 0,9774 & 2.310,0 \\
  \cfg{E2-L} & 50\,\% & 44 & 83 & 0,9731 & 0,9732 & 0,9742 & 2.326,2 \\
\end{longtable}
\endgroup

\subsection{Kiểm thử tự động}
\label{sec:appendix-tests}

Bên cạnh script kiểm tra một lần, dự án có bộ kiểm thử tự động trong \code{tests/test\_m8.py}
được viết cho \code{pytest}. Các bài kiểm thử bám vào những tính chất \emph{bắt buộc} của giao
thức, nên nếu ai đó sửa mã làm sai giao thức thì kiểm thử sẽ thất bại. Nhóm được kiểm tra gồm:

\begin{itemize}[leftmargin=1.2cm,itemsep=2pt]
  \item ánh xạ siêu lớp và ngân sách nhãn khớp bảng~\ref{tab:label-budget};
  \item chuỗi kích thước tensor và số tham số khớp bảng~\ref{tab:arch} và
        \tabref{tab:inference-cost};
  \item stem không có MaxPool, đúng \code{Conv2d(3,64,3,stride=1,padding=1,bias=False)};
  \item \cfg{E1} không có Rotation Head; hai cấu hình ghép cặp dùng cùng trọng số khởi tạo của
        encoder và Superclass Head;
  \item nhãn $k$ sinh ra khớp chính xác \code{torch.rot90} của ảnh gốc, $k$ trải đều cả bốn giá
        trị;
  \item hằng số chuẩn hóa cố định: ảnh đen thành $-1$, ảnh trắng thành $+1$;
  \item phép chia benchmark: kích thước, phân tầng, tập lồng nhau, không rò rỉ sang validation;
  \item \cfg{E1} không tạo DataLoader nào cho nhánh xoay; \cfg{E2-L} chỉ dùng $\DL$; \cfg{E2} dùng
        đủ 45.000 ảnh của $\Dtrain$;
  \item dataset xoay không mang nhãn siêu lớp nên không thể rò rỉ nhãn;
  \item Macro-F1 khớp ví dụ tính tay và quy ước mẫu số bằng $0$ cho kết quả $0$;
  \item bộ phân loại luôn đoán Animal đạt đúng Accuracy $60\,\%$ trên phép chia này;
  \item cấu hình optimizer và scheduler đúng tham số, learning rate về $0$ sau $T_{\max}$ bước;
  \item \textbf{$\lamrot = 0$ vẫn làm thay đổi thống kê BatchNorm} — kiểm chứng trực tiếp lập luận
        ở \secref{sec:exempt-bn};
  \item khởi tạo theo seed là tái lập được;
  \item \textbf{nhánh contrastive}: NT-Xent khớp công thức tính tay và bằng $\log 3$ ở trường
        hợp suy biến; SupCon khớp công thức $L_{\mathrm{out}}$ tính tay và trả $0$ (không
        $\mathrm{NaN}$) khi không có positive pair — đây là hồi quy cho lỗi $-\infty \cdot 0$;
  \item Projection Head không làm thay đổi số tham số suy luận và không nhận gradient từ
        nhiệm vụ phân loại;
  \item dataset contrastive từ $\Dtrain$ \textbf{không} mang nhãn; hai view của cùng một ảnh
        thực sự khác nhau;
  \item ánh xạ \cfg{E3}--\cfg{E6} sang nhánh xoay và loại loss contrastive đúng bảng 5.
\end{itemize}

Kết quả lần chạy gần nhất: \textbf{49 bài kiểm thử ĐẠT, 0 thất bại}
(\code{pytest tests -q}, log lưu ở \code{results/logs/pytest.log}). Ngoài ra
\code{scripts/check\_contrastive.py} chạy kiểm tra độc lập cho nhánh contrastive và trả mã thoát
khác $0$ nếu bất kỳ kiểm tra nào không đạt.

\subsection{Bảng số liệu dạng CSV}
\label{sec:appendix-csv}

Danh sách đầy đủ các chỉ số của từng lượt chạy, ở dạng dễ xử lý bằng máy, được lưu tại
\code{results/runs\_table.csv}. Tệp này chứa: tag lượt chạy, cấu hình, mức nhãn, lần lặp, epoch
được chọn, Macro-F1 và Accuracy trên validation, Macro-F1 và Accuracy trên test, F1 và
precision/recall của từng siêu lớp, thời gian huấn luyện, và số tham số suy luận/huấn luyện.
Ma trận nhầm lẫn đầy đủ của từng lượt nằm trong trường \code{test.confusion} của tệp
\code{results/checkpoints/<tag>.json} tương ứng (không kèm trong kho mã nguồn — xem hộp
\enquote{Lưu ý về gói bàn giao} ở trên).

\subsection{Đường học của toàn bộ lịch sử huấn luyện}
\label{sec:appendix-history}

Mỗi tệp \code{results/checkpoints/<tag>.json} chứa trường \code{history} với đầy đủ 100 bản ghi
epoch, mỗi bản ghi gồm: loss huấn luyện của nhánh có giám sát, loss huấn luyện của nhánh phụ
(bằng \code{null} với \cfg{E1}), tổng loss, loss phân loại trên validation, Macro-F1 validation,
Accuracy validation, learning rate hiện hành và thời gian của epoch. Nhờ vậy, mọi hình trong báo
cáo đều có thể được dựng lại từ dữ liệu thô mà không cần chạy lại mô hình.

\subsection{Ghi nhận về các thiết lập được đề xuất}
\label{sec:appendix-proposed}

Một số hằng số trong báo cáo là \emph{thiết lập được khóa trước để khép kín giao thức}, không
phải giá trị do giảng viên hay sinh viên xác nhận trước khi chạy: seed chia split $2026$, hằng
số chuẩn hóa $(0{,}5;\ 0{,}5;\ 0{,}5)$, quy tắc lấy minibatch cho nhánh phụ, và số lượt đo độ
trễ ($50$ lượt khởi động $+$ $1.000$ lượt đo $\times$ ba phiên). Các giá trị này được chốt
\emph{trước} khi chạy để tránh việc điều chỉnh giao thức sau khi đã thấy kết quả test.
